% problem-000184 6 九硼杂双环壬烷合成
\section*{6. 九硼杂双环壬烷合成}
\textit{下列为分问。}

\subsection*{6-1}
⾦属M的硝酸盐与1,3,5-均苯三甲酸（简写为 $\mathrm { H _ { 3 } B T C }$ ，分⼦量为210.14）按特定⽐例在⼄⼆醇和⽔的混合体系中于 $1 8 0 ~ ^ { \circ } \mathrm { C }$ 下反应12⼩时，得到⼀种具有“孔笼—孔道”结构的⾦属有机⻣架材料Y，密度 $\dot { } \rho _ { } = 0 . 9 6 \mathrm { g }$ $\mathrm { c m } ^ { - 3 } .$ 。单晶X射线衍射分析表明，Y属于⽴⽅晶系，其三维⻣架结构主体由M和 $\mathrm { ( B T C ) ^ { 3 } }$ −组成，呈电中性，结构中带有⽔分⼦。元素分析表明，Y中含M，C，O和H四种元素，且M与C的原⼦⽐为1:6。热重-质谱联合分析结果显⽰，样品在 $1 0 0 { \sim } 2 0 0 \ ^ { \circ } \mathrm { C }$ 区间失重8.21%，对应于化学式中3个⽔分⼦的脱除，⽽主体⻣架结构依然保持；继续加热到 $3 5 0 ~ ^ { \circ } \mathrm { C }$ 失重63.7%，之后⽆明显失重，残渣为⾦属氧化物 $\mathbf { M } \mathrm { O } _ { \circ }$ 通过计算确认M是何种⾦属，写出Y的化学式。

\subsection*{6-2}
将⾦属M加⼊到⾜量的浓硫酸中，微热⽚刻即有⿊⾊物质A⽣成，之后A逐渐转变为灰⽩⾊物质B。该灰⽩⾊物质⽤过量氨⽔充分处理，过滤后，向滤液中通⼊ $\mathrm { S O _ { 2 } }$ ⾄微酸性，⽣成⽩⾊沉淀C（反应1）。元素分析结果，N含量8.65 %，S含量19.6%，H含量2.49 %。测试分析发现，C的结构中负离⼦呈三⻆锥形，磁性测量显⽰抗磁性。C与⾜量 $\mathrm { H _ { 2 } S O _ { 4 } }$ 混合并加热，可⽣成超细粉末态M（反应2）。

\subsection*{6-2}
-1 写出 A，B，C 的化学式。

\subsection*{6-2-2 写出反应1和反应2的⽅程式。}
